\documentclass[12pt]{article}
\usepackage[utf8]{inputenc}
\usepackage[T1]{fontenc}
\usepackage{hyperref}
\title{Tenino Language Resources for K-12 Curriculum Development}
\author{Compiled by Claude Sonnet 5 (claude-sonnet-5) for Andr\'as Kornai}
\date{}
\begin{document}
\maketitle

\section{Introduction}
Tenino (ISO 639-3 tqn) is a Sahaptin variety of the Sahaptian family, associated with the Warm Springs Reservation in Oregon. It is spoken as a first language by older adults only. The search turned up very little material specific to Tenino or Warm Springs. Most of what exists concerns the closely related Yakama (Yakima) Sahaptin dialect, and its fit to Tenino has not been verified. Anyone using it for Tenino should check it with Warm Springs speakers.

\section{Primary Educational Resources}
No Tenino-specific primers, textbooks or school programmes were found in the search results. The University of Washington Sahaptin site may hold teaching materials, but I did not confirm its contents:
\begin{itemize}
\item Sahaptin resource site, University of Washington: \url{https://depts.washington.edu/sahaptin}
\end{itemize}

\section{Linguistic and Reference Materials}
\begin{itemize}
\item Beavert and Hargus, \emph{Ichishkiin Sinwit: Yakama/Yakima Sahaptin Dictionary} (2010), covering the Yakama dialect, not Tenino: \url{https://www.tatteredcover.com/products/9780295989150}
\item A review of the dictionary: \url{https://muse.jhu.edu/article/462418}
\item Overview of Sahaptin on Wikipedia: \url{https://en.wikipedia.org/wiki/Sahaptin_language}
\item Tenino language fact sheet (DICE, University of Missouri): \url{https://dice.missouri.edu/assets/docs/north-america-other/Tenino.pdf}
\item Sharon Hargus, ``Dictionaries and linguists'': \url{https://faculty.washington.edu/sharon/Dictionaries_and_linguists.pdf}
\end{itemize}

\section{Religious and Community Resources}
No Tenino scripture or community-organisation resources were found in the search results. Warm Springs tribal language programmes exist in principle, but I did not find or verify any.

\section{Digital and Audio Resources}
The review of the Yakama dictionary says it includes audio, and these are Yakama recordings, not Tenino. No Tenino-specific audio or digital resources were found.

\section{Curriculum Development Context}
Little verified material exists for Tenino. The language is spoken only by older adults, so documentation with the remaining speakers is the most urgent need. Yakama Sahaptin materials could serve as a starting reference, with adaptation and review by Tenino speakers. Local follow-up with the Warm Springs community is needed.
\end{document}